\section{Conclusion}
\label{sec:discussion}

\paragraph{Interpretation.}
The results support three complementary claims. First, fixed-budget cellular selection improves realism, making both compute and unfiltered performance necessary reporting items. Second, independent analyzers and the random-real-tile control show that spatial agreement is specific to the requested map rather than plausible tissue in general. Third, same-seed dose sweeps recover ordered morphology changes while modeled nuisance factors remain fixed. Analyzer differences keep these claims appropriately narrow: \hovernet{} confirms typed localization, \stardist{} confirms segmentation geometry but not class, and solidity remains sensitive to boundary conventions.

\paragraph{Limitations and Future Work.}
\method{} produces condition-matched samples rather than exact reconstructions; unmeasured properties may still change under shared noise. Global controls and a five-class map compress richer cellular phenotypes. Evaluation is held out but remains within the same broad H\&E domain, motivating external-site, organ-level, pathology-feature, and expert-review studies. Eight-seed selection also trades compute for precision and may behave differently for rare morphologies or sites. Final reporting should include patient/slide confidence intervals and subgroup acceptance rates. Downstream score changes characterize model behavior under synthetic interventions, not patient-level biological causality.

\paragraph{Conclusion.}
\method{} combines dense cellular maps, FiLM morphology controls, same-noise generation, and prediction-independent selection to create fidelity-audited H\&E counterfactuals. Improved realism, cross-analyzer spatial agreement, a random-tile negative control, and monotonic morphology response together support its use as a structured perturbation engine for comparing pathology models and testing explanation faithfulness.
